% Page de garde. Aucune image du système d'information, ici comme ailleurs ; le
% seul logotype admis est celui de l'établissement de formation, et il vient de
% `\logoEcole`, qui compose un cadre de la même hauteur tant que le fichier
% n'existe pas.
%
% Aucune valeur n'est écrite ici : chaque champ vient d'un marqueur.
%
% AUCUN TEXTE FIXE N'ENTOURE UN MARQUEUR QUI PEUT ÊTRE VIDE. Chaque champ dont le
% marqueur peut ne pas être renseigné passe par `\siRenseigne` : l'étiquette
% disparaît avec sa valeur au lieu d'annoncer une phrase qui s'arrête. La règle
% vaut pour les deux noms, pour la date de soutenance, et pour les trois champs de
% contexte. Le seul champ inconditionnel est l'organisme d'accueil, dont le
% marqueur est renseigné dans le fichier des marqueurs lui-même — et il tient la
% table des noms non vide, ce qui évite un tableau sans aucune ligne.

\begin{titlepage}
\centering
\vspace*{1cm}

\logoEcole{2.2cm}

\vspace{0.8cm}

\siRenseigne{\marqueurEtablissementFormation}{%
  {\large \marqueurEtablissementFormation\par}\vspace{0.5cm}%
}
\siRenseigne{\marqueurFiliere}{{\large \marqueurFiliere\par}}

\vfill

\siBrouillon{{\large\fbox{\mentionVersionDeTravail}\par}\vspace{1cm}}

{\Huge\bfseries Rapport de stage d'application\par}
\vspace{1cm}
{\Large Chaîne de données et tableau de bord\\pour un établissement hospitalier\par}

\vfill

\begin{tabular}{rl}
  \siRenseigne{\marqueurAuteur}{Réalisé par & \marqueurAuteur \\[0.4em]}
  \siRenseigne{\marqueurEncadrantProfessionnel}{Encadrant de stage & \marqueurEncadrantProfessionnel \\[0.4em]}
  Organisme d'accueil & \marqueurOrganismeAccueil \\
\end{tabular}

\vfill

\siRenseigne{\marqueurDateSoutenance}{%
  {\large Soutenu le \marqueurDateSoutenance\par}\vspace{0.5cm}%
}
\siRenseigne{\marqueurAnneeUniversitaire}{%
  {\large Année universitaire \marqueurAnneeUniversitaire\par}%
}

\end{titlepage}
